\documentclass[10pt,a4paper]{article}

\usepackage[T1]{fontenc}
\usepackage[utf8]{inputenc}
\usepackage[ngerman]{babel}
\usepackage[a4paper,margin=14mm]{geometry}
\usepackage{lmodern}
\usepackage[table]{xcolor}
\usepackage{tabularx}
\usepackage{array}
\usepackage{microtype}
\usepackage[most]{tcolorbox}

\pagestyle{empty}
\setlength{\parindent}{0pt}
\setlength{\parskip}{0pt}
\renewcommand{\familydefault}{\sfdefault}
\renewcommand{\arraystretch}{1.16}

\definecolor{Navy}{HTML}{17324D}
\definecolor{Blue}{HTML}{3B6EA5}
\definecolor{Sky}{HTML}{EAF3FA}
\definecolor{Mint}{HTML}{E8F5EF}
\definecolor{Gold}{HTML}{F3B33D}
\definecolor{Ink}{HTML}{1E2935}
\definecolor{Muted}{HTML}{617080}
\definecolor{Line}{HTML}{D8E1E8}
\definecolor{CodeBg}{HTML}{F5F7F9}

\color{Ink}
\newcommand{\sectiontitle}[1]{%
  \vspace{2.3mm}{\color{Blue}\bfseries\large #1}\par
  \vspace{0.7mm}{\color{Blue}\rule{\linewidth}{0.7pt}}\par\vspace{1.2mm}}
\newcommand{\schema}[1]{\colorbox{CodeBg}{\parbox{\dimexpr\linewidth-2\fboxsep\relax}{\ttfamily\small #1}}}

\begin{document}
\colorbox{Navy}{\parbox{\dimexpr\linewidth-2\fboxsep\relax}{%
  \vspace{3mm}\color{white}\begin{tabularx}{\linewidth}{@{}Xr@{}}
  {\small\bfseries INFORMATIK · KLASSE 10} & \colorbox{Gold}{\color{Navy}\small\bfseries\strut LÖSUNGEN}\end{tabularx}
  \vspace{2.4mm}{\fontsize{18}{21}\selectfont\bfseries Aufgabe 10 -- 1:1-Beziehungen}\par
  \vspace{0.8mm}{\large Sicherungsblatt}\vspace{2.8mm}}}
\vspace{2.5mm}
\colorbox{Sky}{\parbox{\dimexpr\linewidth-2\fboxsep\relax}{\vspace{1.6mm}
\textcolor{Blue}{\bfseries Ziel:} Du erkennst 1:1-Zuordnungen und erklärst den möglichen Speicherort des Fremdschlüssels.\vspace{1.6mm}}}

\sectiontitle{1 · Klasse und Klassenleiter}
{\small Für dieses Schulbeispiel gilt: Jede Klasse hat genau einen Klassenleiter, und jeder Klassenleiter leitet genau eine Klasse. Die Zuordnung lautet: 10a -- KOH, 10b -- SEI, 10c -- BRA.}\par\vspace{1mm}
{\centering\bfseries\textcolor{Blue}{Klasse \quad 1 \;\rule{32mm}{0.7pt}\; 1 \quad Klassenleiter}\par}
\vspace{1mm}
{\small Ein Fremdschlüssel kann den Verweis in \texttt{klasse} oder in \texttt{klassenleiter} speichern. Schema C mit zwei Verweisen speichert die Verbindung doppelt und kann widersprüchlich werden.}\par
\sectiontitle{2 · Schlüssel und Schloss}
{\small Für diese Schließanlage gilt: Jeder Schlüssel öffnet genau ein Schloss, und zu jedem Schloss gehört genau ein Schlüssel.}\par\vspace{1mm}
{\centering\bfseries\textcolor{Blue}{Schlüssel \quad 1 \;\rule{32mm}{0.7pt}\; 1 \quad Schloss}\par}
\vspace{1.5mm}
\begin{minipage}[t]{0.48\linewidth}
\textcolor{Blue}{\bfseries Möglichkeit A}\par\vspace{0.6mm}
\schema{schluessel (id: int, nummer: int,\par\quad besitzer: varchar(255),\par\quad schloss\_id[schloss]: int)\par
schloss (id: int, ort: varchar(255))}
\end{minipage}\hfill
\begin{minipage}[t]{0.48\linewidth}
\textcolor{Blue}{\bfseries Möglichkeit B}\par\vspace{0.6mm}
\schema{schluessel (id: int, nummer: int,\par\quad besitzer: varchar(255))\par
schloss (id: int, ort: varchar(255),\par\quad schluessel\_id[schluessel]: int)}
\end{minipage}
\par\vspace{1.2mm}
\sectiontitle{3 · Eindeutigkeit}
{\small Ein Fremdschlüssel speichert nur einen Verweis. Damit die Beziehung 1:1 bleibt, darf derselbe Partner nicht mehreren Zeilen zugeordnet werden. Diese Eindeutigkeit gilt in unseren Beispielen als Annahme.}\par\vspace{1mm}
\begin{tcolorbox}[colback=Mint,colframe=Blue,arc=2mm,boxrule=.8pt,left=3mm,right=3mm,top=1.7mm,bottom=1.7mm]
\textbf{Merksatz}\par\smallskip
Bei einer 1:1-Beziehung gehört zu jedem Datensatz der einen Tabelle genau ein Datensatz der anderen Tabelle -- und umgekehrt. Der Fremdschlüssel kann in einer der beiden Tabellen stehen; die Zuordnung muss zusätzlich eindeutig sein.
\end{tcolorbox}

\vfill{\color{Line}\rule{\linewidth}{.5pt}}\par\vspace{1.5mm}
\begin{tabularx}{\linewidth}{@{}Xr@{}}{\small\color{Muted}Klasse 10 · Datenbanken} & {\small\color{Muted}Sicherungsblatt · Aufgabe 10}\end{tabularx}
\end{document}
